\documentclass{article}
\usepackage{amsmath,hyperref}
\begin{document}
\section*{Acquiring the lost aperture distinction before action}
\paragraph{Queries and source.}
Retain the same hold profile (PR 866), including the two possible
left depth/start/duration triples and the possibility that a
permitted left pulse is omitted. Run the selected pinned Open-RMF
velocity functions and publish native mode on the resulting
\emph{modeled} positions. The application must decide whether to
permit the new right pulse at tick 11 with deficit budget 36. The
earlier traversal is safe in all four histories if this request is
held. No fault can withdraw the established pose and scope during
the decision in this fixture. Sample times below index completed
model updates; a qualified sample must arrive by the decision
cut. No ROS network, real sensor, or hardware calibration is
claimed.

\paragraph{Exact profile query versus one new decision.}
The old profile identity is sufficient to answer all the new
prospective aperture queries by its refined index. The particular
tick-11 decision can also depend on whether the left pulse
actually happened. Only the depth-three history \emph{without}
its pulse fails for that right request. The depth-three history
with its pulse and both depth-five histories can permit it.
The resulting certification coarsens by the actual decision,
not by the profile label. The native categorical mode on the
sample after update 3 identifies the old-profile distinction and
allows one safe group, but cannot tell both depth-three cases
apart yet. On the sample after update 4 the same mode decides
the single action exactly. A stale or missing report has no such
credit. All decision arms keep the old traversal, including the
case where they decline the new pulse.

\paragraph{Numerical-only sensor with offset.}
For a fixed run let each numerical observation be the true deficit
plus one unknown additive offset and a bounded per-read error
$\eta_i\in[-e,e]$. The offset is arbitrary but constant within the
run. For competing histories $x,y$ and cuts $T$, they admit a
common transcript iff some offset difference makes
$|x(t)-y(t)+\delta|\le 2e$ for all $t\in T$, equivalently
\[
 \max_{t\in T}(x(t)-y(t))-\min_{t\in T}(x(t)-y(t))\le4e.
\]
This is an interval intersection, not a stochastic independence
assumption. The endpoint implication is checked in Lean.
For complete profile discrimination by update 5, the best pair
has minimum separation span 6, and updates $(1,2,5)$ attain span
8. At $e=7/4$ no pair suffices, whereas those three reads do.
This lower bound includes adaptive two-read schemes: an arbitrary
offset makes any first value, for example zero, possible in
both histories. For the second time chosen on that value, a
paired overlapping transcript can be normalized to start with
exactly that value. Waiting until update 8 permits a successful
two-read profile test at $(3,8)$.
For only the requested right pulse, the relevant separation is
between one unsafe history and all three permitting histories:
updates $(1,4)$ attain span 10 and suffice at $e=7/4$; equality
$e=5/2$ destroys that guarantee.

\paragraph{Operational comparison.}
Replay checks compatible histories of actual biased-read transcripts
on a quarter-unit grid against the unchanged admission core. The
interval argument covers the full bounded real model. The direct
same-evidence interval rule ties. A fresh categorical sample is a
stronger option than numerical-only reconstruction when available
at the right time; there is no general advantage from using more
reads. These are separately indexed observation models and costs,
not measured bandwidth, safety certification, or independent
mission closure.
\end{document}
